\section{Finite optimization with nonnegative Gram entries}
\label{S:finite-positive-gram}

Let $H$ be any real Hilbert space, $\mathcal{I}$ a finite index set, $k_i \in
H$, $\ell_i > 0$ and $b_i \ge 0$. Write $C_{ij} = \langle k_i,k_j\rangle$ and
assume $C_{ij} \ge 0$ for every pair. The Gram matrix is symmetric and positive
semidefinite, since $\sum_{i,j}a_iC_{ij}a_j = \|\sum_i a_i k_i\|^2$. Put $t_+ =
\max(t,0)$ and define
\begin{equation}\label{E:finite-gram-score}
  \mathfrak{F}_C(\Phi)
  =2\sum_{i\in\mathcal{I}}\ell_i
  \sqrt{(\langle\Phi,k_i\rangle-b_i)_+}-\|\Phi\|^2/2,
  \qquad \mathscr{S}(C)=\sup_{\Phi\in H}\mathfrak{F}_C(\Phi).
\end{equation}
Let $\mathcal{J} = \{i:k_i\ne0\} = \{i:C_{ii}>0\}$. Zero vectors contribute zero
because $b_i \ge 0$. Empty index sets, singular Grams and $H = \{0\}$ are
allowed. For nonempty $A \subset \mathcal{J}$ and $w_i > 0$, set
\begin{equation}\label{E:finite-gram-dual}
  D_A^C(w)=\sum_{i\in A}\frac{\ell_i^2}{w_i}
  -\sum_{i\in A}b_iw_i
  +\frac12\sum_{i,j\in A}w_iC_{ij}w_j,
  \qquad d_A^C=\inf_{w_i>0}D_A^C(w).
\end{equation}
Set $d_\varnothing^C = 0$.

\begin{lemma}[Scalar and Hilbert identities]\label{L:gram-scalar-algebra}
  For real $\ell$, $w > 0$ and $t \ge 0$,
  \begin{equation}\label{E:gram-sqrt-dual}
    2\ell\sqrt t\le\ell^2/w+wt.
  \end{equation}
  For $\Phi,z \in H$,
  \begin{equation}\label{E:gram-hilbert-square}
    \langle\Phi,z\rangle-\|\Phi\|^2/2
    =\|z\|^2/2-\|\Phi-z\|^2/2\le\|z\|^2/2.
  \end{equation}
  For finite real $w_i$, $\|\sum_iw_ik_i\|^2 = \sum_{i,j}w_iC_{ij}w_j$.
\end{lemma}

\begin{proof}
  Since $\sqrt w > 0$ and $\sqrt{wt} = \sqrt w\sqrt t$,
  \[
    0\le(\ell/\sqrt w-\sqrt{wt})^2
    =\ell^2/w+wt-2\ell\sqrt t.
  \]
  This includes $t = 0$. If $\ell,t > 0$, equality holds for $w = \ell/\sqrt t$.
  Expanding $\|\Phi-z\|^2$ proves~\eqref{E:gram-hilbert-square}. Bilinearity of
  the real inner product proves the finite Gram identity. These identities use
  neither compactness of a Hilbert ball nor an inverse covariance matrix. This
  proves Lemma~\ref{L:gram-scalar-algebra}.
\end{proof}

\begin{lemma}[Interior dual minimum]\label{L:gram-dual-minimum}
  Each nonempty $A \subset \mathcal{J}$ has a unique positive minimizer $w^A$ of
  $D_A^C$. At this minimizer,
  \begin{equation}\label{E:gram-dual-stationarity}
    \sum_{j\in A}C_{ij}w_j^A-b_i
    =\frac{\ell_i^2}{(w_i^A)^2}>0,\qquad i\in A.
  \end{equation}
\end{lemma}

\begin{proof}
  Put $c_i = C_{ii} > 0$ and $K_A = \sum_{i\in A}b_i^2/c_i$. Nonnegative entries
  and positive weights give
  \[
    D_A^C(w)\ge
    \sum_i\left(\ell_i^2/w_i-b_iw_i+c_iw_i^2/2\right).
  \]
  Completing the square in $c_iw_i^2/4-b_iw_i$ yields
  \begin{equation}\label{E:gram-dual-coercivity}
    D_A^C(w)\ge
    \sum_i\ell_i^2/w_i+\frac14\sum_i c_iw_i^2-K_A.
  \end{equation}
  Thus, no coordinate can approach zero or infinity in a bounded sublevel while
  a different negative term hides its growth. Take $w^0 = (1,\ldots,1)$, $d_0 =
  D_A^C(w^0)$ and $L_A = d_0+K_A$. Equation~\eqref{E:gram-dual-coercivity} at
  $w^0$ gives $L_A \ge \sum_i\ell_i^2+\sum_i c_i/4 > 0$. Every $w$ with
  $D_A^C(w) \le d_0$ satisfies
  \[
    \ell_i^2/L_A\le w_i\le\sqrt{4L_A/c_i}.
  \]
  The sublevel is nonempty and closed inside this compact box, which lies in the
  open positive orthant. A minimizing sequence chosen in the sublevel has a
  convergent subsequence there. Continuity attains its infimum. Outside this
  sublevel the values exceed $d_0$, so the minimum is global on the entire
  positive orthant.

  The function is smooth in that orthant. Its Hessian on a nonzero vector $\xi$
  has value
  \[
    \sum_i\frac{2\ell_i^2}{w_i^3}\xi_i^2
    +\sum_{i,j}\xi_iC_{ij}\xi_j>0.
  \]
  The first term is strictly positive and the second is a squared Hilbert norm.
  Restriction to any nonconstant line segment is strictly convex, proving
  uniqueness. Each coordinate permits both signs of a small perturbation at the
  interior minimum. Differentiation gives $-\ell_i^2/w_i^2-b_i+\sum_jC_{ij}w_j =
  0$. This is~\eqref{E:gram-dual-stationarity}. This proves
  Lemma~\ref{L:gram-dual-minimum}.
\end{proof}

\begin{proposition}[Positive conic maximizer]\label{P:gram-positive-optimizer}
  The finite real value in~\eqref{E:finite-gram-score} satisfies
  \begin{equation}\label{E:gram-active-duality}
    \mathscr{S}(C)=\max_{A\subset\mathcal{J}}d_A^C .
  \end{equation}
  If a maximizing subset is empty, $\Phi = 0$ attains the supremum. If a
  maximizing subset $A$ is nonempty, a maximizer is $\Phi_A = \sum_{i\in A}w_i^A
  k_i$ with strictly positive coefficients. No uniqueness of the primal
  maximizer or maximizing subset is asserted.
\end{proposition}

\begin{proof}
  For arbitrary $\Phi$, let $A(\Phi) = \{i:\langle\Phi,k_i\rangle>b_i\}$. It is
  contained in $\mathcal{J}$. All other terms in the objective vanish, including
  equality at the threshold. If $A(\Phi)$ is empty, $\mathfrak{F}_C(\Phi) =
  -\|\Phi\|^2/2 \le 0$. Otherwise, for any positive $w$ on $A = A(\Phi)$, set
  $t_i = \langle\Phi,k_i\rangle-b_i > 0$ and $z = \sum_{i\in A}w_ik_i$.
  Lemma~\ref{L:gram-scalar-algebra} gives
  \[
    \mathfrak{F}_C(\Phi)
    \le\sum_i\ell_i^2/w_i-\sum_i b_iw_i
    +\langle\Phi,z\rangle-\|\Phi\|^2/2
    \le D_A^C(w).
  \]
  This holds for every positive $w$, so $\mathfrak{F}_C(\Phi) \le d_A^C$. There
  are finitely many subsets and their values are finite by
  Lemma~\ref{L:gram-dual-minimum}. Taking the supremum bounds $\mathscr{S}(C)$
  above by their maximum.

  Conversely, $\mathfrak{F}_C(0) = 0 = d_\varnothing^C$. For nonempty $A$, take
  its dual minimizer. Equation~\eqref{E:gram-dual-stationarity} makes every
  vector in $A$ strictly active at $\Phi_A$, with
  $\sqrt{\langle\Phi_A,k_i\rangle-b_i} = \ell_i/w_i^A$. Both inequalities in the
  preceding display are then equalities for the sum restricted to $A$, since $z
  = \Phi_A$. This restricted objective equals $d_A^C$. All omitted vectors
  contribute nonnegative terms to the full objective, hence
  $\mathfrak{F}_C(\Phi_A) \ge d_A^C$, even if this value is negative. Maximizing
  over $A$, including the empty set, proves \eqref{E:gram-active-duality} and
  the asserted attainment. For a maximizing nonempty subset the upper and lower
  bounds force $\mathfrak{F}_C(\Phi_A) = d_A^C$; no outside inactivity was
  assumed. All-zero vectors give $\mathscr{S}(C) = 0$ directly. This proves
  Proposition~\ref{P:gram-positive-optimizer}.
\end{proof}

\begin{figure}
  \centering
  \begin{tikzpicture}[x=1cm,y=1cm]
    \begin{scope}
      \fill[hamteal!9] (0,0)--(4,.89655)--(4,3)--(2.1875,3)--cycle;
      \draw[ham/arrow] (0,0)--(2.9,.65) node[right] {$k_1$};
      \draw[ham/arrow] (0,0)--(1.75,2.4) node[above] {$k_2$};
      \draw[ham/arrow,hamteal,very thick] (0,0)--(3.0,2.15)
      node[anchor=west] {$\Phi_A=\sum_{i\in A}w_i^Ak_i$};
      \node[ham/note] at (1.8,-.6) {$\langle k_i,k_j\rangle\ge0,\quad w_i^A>0$};
    \end{scope}
    \begin{scope}[xshift=8.3cm]
      \draw[ham/arrow] (0,0)--(2.2,1.1) node[above] {$k_2=2k_1$};
      \draw[ham/arrow,thick] (0,0)--(1.1,.55) node[below] {$k_1$};
      \filldraw[fill=white,draw=hamgold,thick] (0,0) circle (2pt);
      \node[ham/note,text width=3.2cm] at (1.0,2.15)
      {Singular Gram matrices\\are allowed};
      \node[ham/note,text width=3.2cm] at (1,-1.45)
      {A zero vector contributes zero.\\The active set may be empty.};
    \end{scope}
  \end{tikzpicture}
  \caption{Finite positive combinations, not an inverse covariance formula. For
    a nonempty maximizing active set, a maximizer has positive coefficients on
    its selected vectors. The right panel shows why linear independence must not
    be imposed. The drawing is in a finite-dimensional span and does not
    restrict the ambient Hilbert space. Neither the maximizing subset nor the
    maximizer is asserted to be unique. See
    \ProofRef{Proposition}{P:gram-positive-optimizer}.}
  \label{Fig:positive-gram}
\end{figure}


\begin{lemma}[Real Gram realization]\label{L:gram-realization}
  Every finite real symmetric positive semidefinite matrix is the Gram matrix of
  vectors in a finite-dimensional real Hilbert space. When its entries are
  nonnegative, the value $\mathscr{S}(C)$ is independent of its realization.
\end{lemma}

\begin{proof}
  Induct on the size, with the empty matrix realized in $\{0\}$. Diagonal
  entries are nonnegative. If all vanish, testing on $e_i+t e_j$ gives $2tC_{ij}
  \ge 0$ for every real $t$, hence $C = 0$, realized by zero vectors. Otherwise
  permute a positive diagonal $c$ into the first position and write $C =
  \left[\begin{smallmatrix}c&v^{\mathsf{T}}\\v&A\end{smallmatrix}\right]$. The
  Schur complement $B = A-vv^{\mathsf{T}}/c$ is positive semidefinite: test the
  original quadratic form on $(-v^{\mathsf{T}}y/c,y)$. Realize $B$ by vectors
  $u_i$ using the induction hypothesis. Then $k_1 = (\sqrt c,0)$ and $k_i =
  (v_i/\sqrt c,u_i)$ for the other indices realize $C$. Undo the permutation.
  Entries of $B$ need not be nonnegative; this construction requires neither
  nonnegative coordinates nor complete positivity.

  In any realization, $C_{ii} = 0$ means $k_i = 0$, so its entire row vanishes.
  Formula~\eqref{E:gram-active-duality} depends only on $C,\ell,b$ and these
  eligible indices. It therefore proves realization independence whenever the
  entries are nonnegative, which is the scope in which $\mathscr{S}(C)$ is used.
  This proves Lemma~\ref{L:gram-realization}.
\end{proof}

\begin{proposition}[Entrywise Gram comparison]\label{P:gram-entrywise-comparison}
  Let $C,D$ be real symmetric positive semidefinite matrices on the same finite
  index set, with $0 \le C_{ij} \le D_{ij}$. With the same $\ell_i > 0,b_i \ge
  0$,
  \[
    \mathscr{S}(C)\le\mathscr{S}(D).
  \]
  No positive-semidefinite order on $D-C$ is required.
\end{proposition}

\begin{proof}
  Lemma~\ref{L:gram-realization} defines both values. Their eligible sets
  satisfy $\mathcal{J}_C \subset \mathcal{J}_D$. For nonempty $A \subset
  \mathcal{J}_C$ and positive weights,
  \[
    D_A^D(w)-D_A^C(w)
    =\frac12\sum_{i,j\in A}w_i(D_{ij}-C_{ij})w_j\ge0.
  \]
  Taking infima gives $d_A^C \le d_A^D$. The empty values are both zero. Taking
  maxima and using the inclusion of eligible families proves the claim by
  Proposition~\ref{P:gram-positive-optimizer}. This includes the appearance of
  previously zero vectors. This proves
  Proposition~\ref{P:gram-entrywise-comparison}.
\end{proof}

\begin{remark}[Exact checks and scope]\label{R:gram-controls}
  The empty value is essential: for one vector with $C = 1$, $\ell = 1/10,b =
  4$, the nonempty dual has $D(4) = -3199/400 < 0$, whereas the primal objective
  at zero is zero. Global concavity is unavailable: for $k = \ell = b = 1$ the
  objective $2\sqrt{(\phi-1)_+}-\phi^2/2$ decreases on $[0,1]$ and has infinite
  positive right slope at one. For $k = \ell = 1,b = 0$, differentiation on
  $\phi > 0$ gives the maximum $3/2$ at $\phi = 1$; on $\phi \le 0$ the
  objective is nonpositive. Doubling the cost changes the positive critical
  point to $2^{-2/3}$ and the maximum to $3\,2^{-4/3}$. Finally, $C = I_2$, $D =
  \left[\begin{smallmatrix}1&1/2\\1/2&1\end{smallmatrix}\right]$ are nonnegative
  positive semidefinite Grams but $D-C$ has eigenvalues $1/2,-1/2$. These are
  finite analytic checks, not a covariance-order theorem for physical waves.

  The scalar, finite Gram, conditional active upper and half-quadratic entrywise
  deductions have existing local Lean certificates at their exact hypotheses.
  Dual coercivity, calculus, attainment, realization and the infimum/maximum
  argument above remain analytic proofs. No new kernel fact or formalization of
  the physical application follows merely from these certificates.
\end{remark}
